\chapter{Messverfahren}\label{ch:messverfahren}

% [Cross-Ref Konsolidierte Doku, K-I.2 2026-05-15]
% Plattform-Modell + Hardware-Klassen (Termin 6 + Rang-3 P31-P33): docs/architektur/04_konzepte_saeule_b.md
% Cache-Strategy 29 Familien: docs/architektur/09_taxonomien.md
% Mess-Pipeline Sessions V19-V22: docs/sessions/20260514-4500-habich-termin-8-briefing-refresh-v19-v31.md
% Termin 3 (Benchmark + Datensatzplan, YCSB): docs/termine_konsolidiert/03_termin_3.md

\section{Drei Pflicht-Messreihen}

Die Architektur sieht drei Messreihen vor, die alle vom
\texttt{messung\_driver} (Diplomarbeit/Code/) ueber die
\texttt{ExperimentDriver}-Library (cache-engine) ausgefuehrt werden.
Die Konfiguration erfolgt entweder als drei feste Reihen
(klassischer Modus) oder ueber eine externe
\texttt{messreihen.xml}-Datei (V11.5 Standard-Template).

\subsection{Messreihe A: PRT-ART vs. Stand der Technik}

Messreihe~A vergleicht PRT-ART als Pruefling gegen alle in der
\texttt{cache\_engine/algorithm\_profiles/sota/} dokumentierten SOTA-
Algorithmen. Zwei Modi sind verfuegbar:

\begin{description}
    \item[A\_defined] Schnelle Vorab-Verifikation gegen die 8~Rang-1-
          Profile (art, hot, masstree, coco\_trie, start, b2tree,
          wormhole, surf). Mode \texttt{defined} mit
          \texttt{sota\_profile\_filter} in
          \texttt{ExperimentDriverOptions}.
    \item[A\_full] Vollstaendige Messung gegen alle 30~SOTA-Profile
          (Rang-1 + Rang-2/3 fuer P11--P32). Mode \texttt{full} ist
          Default gemaess Habich-Direktive 2026-05-14
          (``Messung immer so vollstaendig wie moeglich'').
\end{description}

\subsection{Messreihe B: Cache-Engine-Permutationen}

Reine Vergleichsbasis ``Stand der Technik'' --- alle SOTA-Profile ohne
PRT-ART-Pruefling. Im aktuellen Skelett identisch zu A\_full;
zukuenftige Erweiterung kann zusaetzlich
\texttt{cache\_engine\_permutations.xml}-Achsen einbeziehen.

\subsection{Messreihe C: Merge alt/neu}

Konkrete Merge-Punkte zwischen PRT-ART-Bausteinen und cache-engine-
Permutationen:

\begin{itemize}
    \item PRT-ART OLC + tcmalloc-Allocator + ART-Page
    \item PRT-ART 4+2-Pool + HOT-Page
    \item PRT-ART PathPrefetch + B$^2$-Tree-Page
    \item PRT-ART MultiLevelLayout + Wormhole-Page
\end{itemize}

\section{Profile-aware Workload-Routing}

\texttt{ExperimentDriver::phase5\_run\_workload} (V11.2) leitet pro
Permutations-Modul den YCSB-Workload aus dem \texttt{traversal}-Tag
des Profils ab:

\begin{table}[h]
\centering
\begin{tabular}{lll}
\toprule
\textbf{Traversal-Tag} & \textbf{YCSB-Workload} & \textbf{Charakteristik} \\
\midrule
\texttt{RANGE\_SCAN} / \texttt{RANGE\_HOT\_PATH} & YCSB-E & Range-Queries \\
\texttt{HOT\_PATH} / \texttt{PRTART\_HOT\_PATH} & YCSB-A & 50/50 Read/Update Zipfian \\
\texttt{STANDARD} (Default) & YCSB-C & Read-only \\
\bottomrule
\end{tabular}
\caption{Profile-aware Workload-Routing (V11.2)}
\label{tab:workload-routing}
\end{table}

Module ohne zugehoeriges Profil verwenden den globalen Default-Workload
aus \texttt{WorkloadOptions::workload}.

\section{ExperimentDriver Phase 1--7}

Der \texttt{ExperimentDriver} (V8.7 Library + V10.5 Profile-Auto-Pickup)
durchlaeuft sieben Phasen:

\begin{enumerate}
    \item \textbf{Enumerate} --- XML-Configs parsen + Permutations-
          Liste erzeugen
    \item \textbf{Codegen} --- pro Permutation eine
          \texttt{comdare\_perm\_<fp>.cpp} (V9.3 zusaetzlich
          \texttt{generate\_module\_from\_profile} fuer 30~SOTA-Profile)
    \item \textbf{Compile} --- cmake configure + cmake \texttt{--build}
    \item \textbf{Load} --- LoadLibrary/dlopen aller
          \texttt{.dll}/\texttt{.so}-Module
    \item \textbf{Execute} --- pro Modul \texttt{create\_instance} +
          \texttt{run\_workload} (V11.2 Profile-aware)
    \item \textbf{Measure} --- binary measurement-records aggregieren
    \item \textbf{Persist} --- \texttt{measurements.csv} +
          \texttt{measurements.json}
\end{enumerate}

\textbf{Opt-in-Phasen:}

\begin{itemize}
    \item \texttt{phase3\_hot\_compile\_missing} (V13.2): falls
          \texttt{enable\_runtime\_codegen=true} und ein
          vorkompiliertes Modul fehlt, wird per \texttt{cmake/cl}-Aufruf
          hot-kompiliert.
    \item \texttt{phase4b\_functional\_tests} (V13.3): falls
          \texttt{enable\_functional\_tests=true}, wird pro Modul
          die ABI-Vertrags-Konformitaet (\texttt{create\_instance} non-null,
          \texttt{run\_workload} mit empty WorkloadDescriptor schreibt
          \texttt{rec.version == COMDARE\_ABI\_VERSION},
          \texttt{destroy\_instance} sicher) verifiziert.
\end{itemize}

\section{Mess-Modus: COMDARE\_EXPERIMENT\_MODE}

Der Mess-Pfad in der \texttt{ExecutionEngine} (V8.5) ist nur dann aktiv,
wenn die cache-engine mit \texttt{-DCOMDARE\_EXPERIMENT\_MODE=ON}
kompiliert wurde. Im Production-Pfad (Default OFF) ist der
\texttt{ResultAggregator} kein Member der \texttt{ExecutionEngine}, was
Production-Code von jeglichem Mess-Overhead befreit. Der
Diplomarbeits-\texttt{messung\_driver} setzt das Flag explizit beim
Sub-Build (V8.12).

\section{Permutations-Explosion durch 14 Achsen}\label{sec:permexplosion}

% Quelle: docs/bausteine/05_flag_system.md §9 (N.10 Phase 2026-05-18)
% Quelle: docs/bausteine/07_bausteine_matrix_N_erweitert.md §1-§6

Mit der N-Phasen-Erweiterung der Bausteine-Matrix von 11 auf
14~Achsen + ca.~30 Sub-Achsen (siehe Kapitel~\ref{sec:14-achsen-substruktur})
wird der Permutations-Raum substantiell groesser.

\subsection{Quantitative Schaetzung}

\begin{itemize}
    \item \textbf{Vor der N-Phase} ($\sim$11~Achsen, $\sim$85 Bausteine):
    geschaetzt $\sim$5,5~Milliarden Permutationen.
    \item \textbf{Nach der N-Phase} (14~Achsen + Sub-Achsen, ca.~120 Bausteine):
    geschaetzt \textbf{$\sim$100~Milliarden+ Permutationen}.
\end{itemize}

Im PermutationFlags-Bitfeld wachsen entsprechend die Banks:
9~Banks $\rightarrow$ 14~Banks, $\sim$50~Bit $\rightarrow$
82~Bit Permutations-ID-Breite.

\subsection{Konsequenzen fuer die drei Messreihen}

Ein vollstaendiger Permutations-Sweep ueber 100~Mrd Konfigurationen ist
selbst auf dem ZIH-Cluster nicht erschoepfend laufbar. Drei
Mechanismen daempfen die Explosion:

\paragraph{(1) Profile-Filter via \texttt{<expected\_workload>}.}
Seit V19 (\texttt{algorithm\_profiles/}-XML-Schema) tagged jedes
SOTA-Profil seinen erwarteten Workload (z.\,B.\ YCSB\_A fuer Update-heavy,
YCSB\_C fuer Read-only). Phase~5 des \texttt{ExperimentDriver} routet
pro Permutation nur die kompatiblen Workloads -- nicht $6\times$ alles.

\paragraph{(2) \texttt{MessreihenMode::Defined} fuer Manuskript-Plots.}
Standard-Modus liest \texttt{messreihen.xml} und fuehrt nur die
\emph{deklarierten} Tupel aus. Die drei Pflicht-Messreihen
(\texttt{config\_a\_prt\_art\_vs\_sota.xml},
\texttt{config\_b\_cache\_engine\_perms.xml},
\texttt{config\_c\_merge\_alt\_neu.xml}) deklarieren typischerweise
ca.~100--1000 Tupel pro Reihe. Diese Mess-Daten landen direkt im
Manuskript-Anhang.

\paragraph{(3) \texttt{MessreihenMode::Full} nur fuer ZIH-Cluster.}
Der Full-Modus iteriert systematisch ueber den vollen Permutations-Raum
unter Beachtung der Achsen-Kompatibilitaet (z.\,B.\ AVX-512 nur auf
Plattformen mit AVX-512-Support). Realistischer Sub-Modus ist
\texttt{Full-Sampled} (z.\,B.\ jede 1000.\ Permutation), der weiterhin
sinnvolle Stichprobe der ca.~100~Mrd Permutationen liefert. Diese
Laeufe finden auf den drei TIER-3-HABICH-Plattformen statt:
ZIH-Talos-OS-K8s-Pods auf Ryzen~9~9950X3D, Intel~i9~14900KS und
P31-HBM-Hardware (vgl.\ Kapitel~\ref{sec:hw-sched-pro-paper}).

\subsection{Hardware-Pflicht zur Achsen-Reduktion}

Nicht alle 14~Achsen sind auf allen Plattformen variierbar. Beispielsweise:

\begin{itemize}
    \item Achse~12.1 SIMD-Family: AVX-512 ist nur auf
    Server-Plattformen verfuegbar (i9 14900KS hat keine
    AVX-512-Encoding-Unit).
    \item Achse~12.2 Cache-Level-Targeting: HBM-aware ist nur auf
    HBM-Hardware sinnvoll (P31-Rang-3).
    \item Achse~13.2 SIMD-Worker-Count-Limit: hardware-limitiert
    durch die Anzahl physischer SIMD-Einheiten (typisch~2 von $N$
    Cores), nicht durch Software.
\end{itemize}

Der \texttt{IPlatformProbe} (Phase~1 DISCOVER der Builder-Pipeline)
liefert daher pro Plattform eine reduzierte \emph{lauffaehige
Permutations-Achsen-Menge}, die den Permutations-Raum von vornherein
auf das auf dem Host machbare einschraenkt.
